\section{Completion of the mixing bound}\label{sec:completion}

Return to $G=\Sym(V)$ for $n=2^d$. We now pass from the operator estimate
to the full permutation law.
Use the Fourier convention from Section~\ref{sec:lift}, and write
$D_\rho$ for the dimension of an irreducible unitary representation
$\rho$. Our convolution convention gives
$\widehat{\mu*\nu}(\rho)=\widehat\mu(\rho)\widehat\nu(\rho)$.
The following Plancherel bound is the classical Fourier method used by
Diaconis and Shahshahani~\cite{DiaconisShahshahani1981}. We keep the
full matrix form because the shuffle law need not be central, and we
record the normalization explicitly.

\begin{lemma}[The full-group variation bound]\label{lem:plancherel}
For any probability measure $\nu$ on $G$,
\begin{equation}\label{eq:plancherel-tv}
 4\|\nu-U\|_{\TV}^2
 \le\sum_{\rho\ne\mathrm{triv}}D_\rho
       \|\widehat\nu(\rho)\|_{\HS}^2,
\end{equation}
where the sum contains one representative of each nontrivial irreducible
equivalence class.
\end{lemma}
\begin{proof}
Cauchy--Schwarz bounds the left side by
$|G|\sum_g|\nu(g)-|G|^{-1}|^2$.
For any real function $a$ on $G$, expansion and the regular-character
identity give
\begin{align*}
 \sum_\rho D_\rho\Bigl\|\sum_g a(g)\rho(g)\Bigr\|_{\HS}^2
 &=\sum_{g,h}a(g)a(h)
   \sum_\rho D_\rho\tr\bigl(\rho(h)^*\rho(g)\bigr)\\
 &=|G|\sum_g a(g)^2.
\end{align*}
Apply this to $a(g)=\nu(g)-|G|^{-1}$. The trivial coefficient vanishes;
uniform averaging in every nontrivial irreducible is zero. These facts
give \eqref{eq:plancherel-tv}.
\end{proof}

\begin{lemma}[Fair sign]\label{lem:fair-sign}
For every $d\ge1$ and every integer $t\ge1$,
$\widehat{q_d^{*t}}(\sgn)=0$. If a probability measure $\nu$ is
$\varepsilon$ away from $q_d^{*t}$ in total variation, then
$|\widehat\nu(\sgn)|\le2\varepsilon$. If instead
$0\le\xi\le q_d^{*t}$ has mass $m$, then
$|\widehat\xi(\sgn)|\le1-m$.
\end{lemma}
\begin{proof}
Condition on every switch bit except one in the final physical shuffle.
Changing that bit composes the output permutation with a transposition,
so its two possible signs are opposite and equally likely. The
expectation of sign is therefore zero. The two perturbation estimates
follow by summing against a function of absolute value one; their
respective $\ell^1$ differences are $2\varepsilon$ and $1-m$.
\end{proof}

\begin{proof}[Proof of Theorem~\ref{thm:main}]
Fix
$k=64$ and let $d\ge6$ grow. Put $T=1025d$ and
$\mu=q_d^{*T}$. Proposition~\ref{prop:marginals} supplies a partition
of the labels into $k$ sets of size $m=n/k$ for which the sum of the
complementary marginal errors tends to zero. Lemma~\ref{lem:trim}
therefore applies with their sum at most $1/4$ for all sufficiently large
$d$. It supplies a probability measure $\bar\mu$ such that
\begin{equation}\label{eq:final-trim}
 \|\bar\mu-\mu\|_{\TV}=o(1),\qquad
 \bar\mu(gH_a)\le4\frac{|H_a|}{|G|}\quad(g\in G,\ 1\le a\le k).
\end{equation}
Here $H_a$ permutes the $a$th label block and fixes all other labels.
Corollary~\ref{prop:lift} gives,
simultaneously for every irreducible $\rho$ of dimension $D$,
\begin{equation}\label{eq:final-op}
 \|\widehat{\bar\mu}(\rho)\|_{\op}
 \le A D^{-23/64},\qquad A=\sqrt{4\cdot64\,C_0},
\end{equation}
where $C_0$ is the absolute constant in Lemma~\ref{lem:degrees}.

Take $\ell=32$, a fixed number independent of $d$. Operator-norm
submultiplicativity and $\|B\|_{\HS}\le\sqrt D\|B\|_{\op}$ show that
the contribution of the nontrivial, nonsign representations to the right
side of \eqref{eq:plancherel-tv}, for $\nu=\bar\mu^{*\ell}$, is at most
\begin{align}
 \sum_{\rho\ne\mathrm{triv},\sgn}
       D\|\widehat{\bar\mu}(\rho)^\ell\|_{\HS}^2
 &\le A^{64}\sum_{\rho\ne\mathrm{triv},\sgn}
       D^{2-64(23/64)}\notag\\
 &=A^{64}\sum_{\rho\ne\mathrm{triv},\sgn}D^{-21}=o(1).
 \label{eq:final-degree-sum}
\end{align}
The last conclusion follows from Lemma~\ref{lem:degrees}, since
$D^{-21}\le D^{-16}$ for $D\ge1$. Neither diagonalizability nor normality
of the Fourier matrices is needed here.

The sign representation has dimension one, so
\eqref{eq:final-op} alone gives no decay. Lemma~\ref{lem:fair-sign}
shows that $\widehat\mu(\sgn)=0$ and therefore
\[
 |\widehat{\bar\mu}(\sgn)|
 \le2\|\bar\mu-\mu\|_{\TV}=o(1).
\]
Its contribution after $\ell$ factors is therefore $o(1)$ as well.
Together with Lemma~\ref{lem:plancherel} and
\eqref{eq:final-degree-sum}, this proves
$\|\bar\mu^{*32}-U\|_{\TV}=o(1)$.

Convolution by a probability measure contracts total variation.
Replacing the factors one at a time and using \eqref{eq:final-trim} gives
\[
 \|\mu^{*32}-U\|_{\TV}
 \le32\|\mu-\bar\mu\|_{\TV}
       +\|\bar\mu^{*32}-U\|_{\TV}=o(1).
\]
Successive independent $T$-step blocks of the original shuffle have law
$\mu$, so $\mu^{*32}=q_d^{*32800d}$. Equation~\eqref{eq:worst-state}
makes this conclusion uniform over the initial ordering. In particular
the distance is at most $1/4$ for all sufficiently large $d$, as required.
\end{proof}

The constants were chosen to leave room in two simple estimates: the
hidden-card contraction must overcome the number of sequential marginal
comparisons, and the fixed convolution count must overcome the Fourier
dimension weights. No optimal leading constant or cutoff statement is
asserted. Together with the support lower bound, the result gives
mixing order $\Theta(d)$.

\paragraph{Further partial-permutation laws.}
Theorem~\ref{thm:main} is complete with the frame, omitted-block, orbit,
and tableau arguments above. The remaining sections compare estimates for
each fixed list with estimates averaged over lists. They then examine what
can be retained after stopping or conditioning on trajectory events and
develop Fourier and character transfers for those forms of partial
information.
